\documentclass[aspectratio=169, 11pt]{beamer}
% ============================================================================
% Preambulo compartido para presentaciones Beamer
% Estadistica 2026-II - Pontificia Universidad Javeriana
% Incluir con: % ============================================================================
% Preambulo compartido para presentaciones Beamer
% Estadistica 2026-II - Pontificia Universidad Javeriana
% Incluir con: % ============================================================================
% Preambulo compartido para presentaciones Beamer
% Estadistica 2026-II - Pontificia Universidad Javeriana
% Incluir con: \input{../preambulo_beamer.tex} despues de \documentclass
% ============================================================================

\usepackage[T1]{fontenc}
\usepackage[utf8]{inputenc}
\usepackage[spanish]{babel}
\usepackage{amsmath, amssymb, amsfonts}
\usepackage{booktabs}
\usepackage{graphicx}
\usepackage{tikz}
\usetikzlibrary{shapes.geometric, positioning, arrows.meta, calc}
\usepackage{pgfplots}
\pgfplotsset{compat=1.18}
\usepackage{listings}
\usepackage{xcolor}
\usepackage{hyperref}
\usepackage{multicol}

% --- Tema Beamer (minimalista) ---
\usetheme{default}
\usecolortheme{default}
\usefonttheme{default}
\setbeamertemplate{navigation symbols}{}
\setbeamertemplate{caption}[numbered]
\setbeamertemplate{footline}{%
  \hspace{1em}{\tiny\url{https://github.com/polanco-jaime/Curso-Estadistica-2026-3}}\hfill\usebeamerfont{footline}\insertframenumber/\inserttotalframenumber\hspace{1em}\vspace{0.5em}%
}
\setbeamertemplate{itemize items}[circle]
\setbeamertemplate{enumerate items}[default]

% --- Colores institucionales (sutiles) ---
\definecolor{javeriana}{RGB}{0, 62, 105}
\definecolor{javerianalight}{RGB}{0, 105, 170}
\definecolor{codigobg}{RGB}{245, 245, 245}
\definecolor{codigofg}{RGB}{40, 40, 40}
\definecolor{comentario}{RGB}{100, 100, 100}
\definecolor{keyword}{RGB}{0, 100, 150}
\definecolor{string}{RGB}{180, 60, 30}

\setbeamercolor{title}{fg=javeriana}
\setbeamercolor{frametitle}{fg=javeriana}
\setbeamercolor{structure}{fg=javeriana}
\setbeamercolor{block title}{fg=white, bg=javeriana!75}
\setbeamercolor{block body}{bg=javeriana!5}
\setbeamercolor{block title alerted}{fg=white, bg=red!70!black}
\setbeamercolor{block body alerted}{bg=red!5}
\setbeamercolor{block title example}{fg=white, bg=green!40!black}
\setbeamercolor{block body example}{bg=green!5}
\setbeamercolor{item}{fg=javeriana}

\setbeamerfont{title}{series=\bfseries, size=\Large}
\setbeamerfont{frametitle}{series=\bfseries}

% --- Configuracion de listings para R ---
\lstset{
  language=R,
  basicstyle=\ttfamily\footnotesize,
  backgroundcolor=\color{codigobg},
  commentstyle=\color{comentario}\itshape,
  keywordstyle=\color{keyword}\bfseries,
  stringstyle=\color{string},
  numbers=none,
  frame=single,
  rulecolor=\color{javeriana!30},
  breaklines=true,
  showstringspaces=false,
  columns=flexible,
  morekeywords={library, ggplot, aes, geom_histogram, geom_boxplot,
    geom_point, geom_smooth, geom_density, geom_bar, geom_line,
    geom_vline, geom_hline, geom_errorbar, labs, theme_minimal,
    facet_wrap, scale_fill_manual, scale_color_manual,
    summarise, group_by, filter, mutate, select, arrange,
    read.csv, head, str, summary, mean, median, sd, var, cor,
    lm, t.test, chisq.test, wilcox.test, kruskal.test, aov,
    pnorm, qnorm, dnorm, rnorm, qt, pt, qchisq, pchisq, qf, pf,
    confint, coeftest, vcovHC, bptest, vif, cooks.distance,
    TukeyHSD, table, prop.table}
}

% --- Comandos personalizados ---
\newcommand{\R}{\texttt{R}}
\newcommand{\code}[1]{\texttt{\footnotesize #1}}
\newcommand{\variable}[1]{\texttt{#1}}
\newcommand{\paquete}[1]{\texttt{#1}}

% --- Informacion del curso ---
\institute{Pontificia Universidad Javeriana\\
  Facultad de Ciencias Econ\'omicas y Administrativas}
\date{2026-II}
 despues de \documentclass
% ============================================================================

\usepackage[T1]{fontenc}
\usepackage[utf8]{inputenc}
\usepackage[spanish]{babel}
\usepackage{amsmath, amssymb, amsfonts}
\usepackage{booktabs}
\usepackage{graphicx}
\usepackage{tikz}
\usetikzlibrary{shapes.geometric, positioning, arrows.meta, calc}
\usepackage{pgfplots}
\pgfplotsset{compat=1.18}
\usepackage{listings}
\usepackage{xcolor}
\usepackage{hyperref}
\usepackage{multicol}

% --- Tema Beamer (minimalista) ---
\usetheme{default}
\usecolortheme{default}
\usefonttheme{default}
\setbeamertemplate{navigation symbols}{}
\setbeamertemplate{caption}[numbered]
\setbeamertemplate{footline}{%
  \hspace{1em}{\tiny\url{https://github.com/polanco-jaime/Curso-Estadistica-2026-3}}\hfill\usebeamerfont{footline}\insertframenumber/\inserttotalframenumber\hspace{1em}\vspace{0.5em}%
}
\setbeamertemplate{itemize items}[circle]
\setbeamertemplate{enumerate items}[default]

% --- Colores institucionales (sutiles) ---
\definecolor{javeriana}{RGB}{0, 62, 105}
\definecolor{javerianalight}{RGB}{0, 105, 170}
\definecolor{codigobg}{RGB}{245, 245, 245}
\definecolor{codigofg}{RGB}{40, 40, 40}
\definecolor{comentario}{RGB}{100, 100, 100}
\definecolor{keyword}{RGB}{0, 100, 150}
\definecolor{string}{RGB}{180, 60, 30}

\setbeamercolor{title}{fg=javeriana}
\setbeamercolor{frametitle}{fg=javeriana}
\setbeamercolor{structure}{fg=javeriana}
\setbeamercolor{block title}{fg=white, bg=javeriana!75}
\setbeamercolor{block body}{bg=javeriana!5}
\setbeamercolor{block title alerted}{fg=white, bg=red!70!black}
\setbeamercolor{block body alerted}{bg=red!5}
\setbeamercolor{block title example}{fg=white, bg=green!40!black}
\setbeamercolor{block body example}{bg=green!5}
\setbeamercolor{item}{fg=javeriana}

\setbeamerfont{title}{series=\bfseries, size=\Large}
\setbeamerfont{frametitle}{series=\bfseries}

% --- Configuracion de listings para R ---
\lstset{
  language=R,
  basicstyle=\ttfamily\footnotesize,
  backgroundcolor=\color{codigobg},
  commentstyle=\color{comentario}\itshape,
  keywordstyle=\color{keyword}\bfseries,
  stringstyle=\color{string},
  numbers=none,
  frame=single,
  rulecolor=\color{javeriana!30},
  breaklines=true,
  showstringspaces=false,
  columns=flexible,
  morekeywords={library, ggplot, aes, geom_histogram, geom_boxplot,
    geom_point, geom_smooth, geom_density, geom_bar, geom_line,
    geom_vline, geom_hline, geom_errorbar, labs, theme_minimal,
    facet_wrap, scale_fill_manual, scale_color_manual,
    summarise, group_by, filter, mutate, select, arrange,
    read.csv, head, str, summary, mean, median, sd, var, cor,
    lm, t.test, chisq.test, wilcox.test, kruskal.test, aov,
    pnorm, qnorm, dnorm, rnorm, qt, pt, qchisq, pchisq, qf, pf,
    confint, coeftest, vcovHC, bptest, vif, cooks.distance,
    TukeyHSD, table, prop.table}
}

% --- Comandos personalizados ---
\newcommand{\R}{\texttt{R}}
\newcommand{\code}[1]{\texttt{\footnotesize #1}}
\newcommand{\variable}[1]{\texttt{#1}}
\newcommand{\paquete}[1]{\texttt{#1}}

% --- Informacion del curso ---
\institute{Pontificia Universidad Javeriana\\
  Facultad de Ciencias Econ\'omicas y Administrativas}
\date{2026-II}
 despues de \documentclass
% ============================================================================

\usepackage[T1]{fontenc}
\usepackage[utf8]{inputenc}
\usepackage[spanish]{babel}
\usepackage{amsmath, amssymb, amsfonts}
\usepackage{booktabs}
\usepackage{graphicx}
\usepackage{tikz}
\usetikzlibrary{shapes.geometric, positioning, arrows.meta, calc}
\usepackage{pgfplots}
\pgfplotsset{compat=1.18}
\usepackage{listings}
\usepackage{xcolor}
\usepackage{hyperref}
\usepackage{multicol}

% --- Tema Beamer (minimalista) ---
\usetheme{default}
\usecolortheme{default}
\usefonttheme{default}
\setbeamertemplate{navigation symbols}{}
\setbeamertemplate{caption}[numbered]
\setbeamertemplate{footline}{%
  \hspace{1em}{\tiny\url{https://github.com/polanco-jaime/Curso-Estadistica-2026-3}}\hfill\usebeamerfont{footline}\insertframenumber/\inserttotalframenumber\hspace{1em}\vspace{0.5em}%
}
\setbeamertemplate{itemize items}[circle]
\setbeamertemplate{enumerate items}[default]

% --- Colores institucionales (sutiles) ---
\definecolor{javeriana}{RGB}{0, 62, 105}
\definecolor{javerianalight}{RGB}{0, 105, 170}
\definecolor{codigobg}{RGB}{245, 245, 245}
\definecolor{codigofg}{RGB}{40, 40, 40}
\definecolor{comentario}{RGB}{100, 100, 100}
\definecolor{keyword}{RGB}{0, 100, 150}
\definecolor{string}{RGB}{180, 60, 30}

\setbeamercolor{title}{fg=javeriana}
\setbeamercolor{frametitle}{fg=javeriana}
\setbeamercolor{structure}{fg=javeriana}
\setbeamercolor{block title}{fg=white, bg=javeriana!75}
\setbeamercolor{block body}{bg=javeriana!5}
\setbeamercolor{block title alerted}{fg=white, bg=red!70!black}
\setbeamercolor{block body alerted}{bg=red!5}
\setbeamercolor{block title example}{fg=white, bg=green!40!black}
\setbeamercolor{block body example}{bg=green!5}
\setbeamercolor{item}{fg=javeriana}

\setbeamerfont{title}{series=\bfseries, size=\Large}
\setbeamerfont{frametitle}{series=\bfseries}

% --- Configuracion de listings para R ---
\lstset{
  language=R,
  basicstyle=\ttfamily\footnotesize,
  backgroundcolor=\color{codigobg},
  commentstyle=\color{comentario}\itshape,
  keywordstyle=\color{keyword}\bfseries,
  stringstyle=\color{string},
  numbers=none,
  frame=single,
  rulecolor=\color{javeriana!30},
  breaklines=true,
  showstringspaces=false,
  columns=flexible,
  morekeywords={library, ggplot, aes, geom_histogram, geom_boxplot,
    geom_point, geom_smooth, geom_density, geom_bar, geom_line,
    geom_vline, geom_hline, geom_errorbar, labs, theme_minimal,
    facet_wrap, scale_fill_manual, scale_color_manual,
    summarise, group_by, filter, mutate, select, arrange,
    read.csv, head, str, summary, mean, median, sd, var, cor,
    lm, t.test, chisq.test, wilcox.test, kruskal.test, aov,
    pnorm, qnorm, dnorm, rnorm, qt, pt, qchisq, pchisq, qf, pf,
    confint, coeftest, vcovHC, bptest, vif, cooks.distance,
    TukeyHSD, table, prop.table}
}

% --- Comandos personalizados ---
\newcommand{\R}{\texttt{R}}
\newcommand{\code}[1]{\texttt{\footnotesize #1}}
\newcommand{\variable}[1]{\texttt{#1}}
\newcommand{\paquete}[1]{\texttt{#1}}

% --- Informacion del curso ---
\institute{Pontificia Universidad Javeriana\\
  Facultad de Ciencias Econ\'omicas y Administrativas}
\date{2026-II}

\title{Sesión 1: Estadística Descriptiva Numérica}
\subtitle{Introducción al curso y medidas numéricas}
\author{Prof.\ Jaime Polanco Jim\'enez}
\date{Vie 31 jul 2026}
\begin{document}

\begin{frame}
\titlepage
\end{frame}

\begin{frame}{Agenda}
\tableofcontents
\end{frame}

% ============================================================================
\section{Introducción al curso}
% ============================================================================

\begin{frame}{¿Qué es la estadística?}
\begin{block}{Definición}
La estadística es la ciencia de recolectar, organizar, analizar e interpretar datos para tomar decisiones informadas bajo incertidumbre.
\end{block}

\vspace{0.5em}
\begin{columns}[T]
\column{0.48\textwidth}
\textbf{Estadística Descriptiva}
\begin{itemize}
\item Resumir y visualizar datos
\item Medidas numéricas (media, mediana, desviación estándar)
\item Gráficos (histogramas, box plots, scatter plots)
\item \alert{Describe lo que ya tenemos}
\end{itemize}

\column{0.48\textwidth}
\textbf{Estadística Inferencial}
\begin{itemize}
\item Generalizar desde una muestra a la población
\item Intervalos de confianza
\item Pruebas de hipótesis
\item Regresión y predicción
\item \alert{Infiere sobre lo que no observamos}
\end{itemize}
\end{columns}
\end{frame}

\begin{frame}{Datos ICFES: Saber 11 y Saber Pro}
\begin{block}{¿Qué es el ICFES?}
Instituto Colombiano para la Evaluación de la Educación. Realiza exámenes estandarizados para evaluar la calidad educativa en Colombia.
\end{block}

\begin{itemize}
\item \textbf{Saber 11}: examen al finalizar la educación media (bachillerato)
  \begin{itemize}
  \item Módulos: Lectura Crítica, Matemáticas, Sociales, Ciencias Naturales, Inglés
  \item Usado para admisión universitaria
  \end{itemize}
\item \textbf{Saber Pro}: examen al finalizar la educación superior (último año de pregrado)
  \begin{itemize}
  \item Módulos genéricos: Lectura Crítica, Razonamiento Cuantitativo, Competencias Ciudadanas, Comunicación Escrita, Inglés
  \item Módulos específicos por programa (ej: Gestión Financiera para Ciencias Econ\'omicas)
  \end{itemize}
\end{itemize}
\end{frame}

\begin{frame}{Datos ICFES: enfoque del curso}
\alert{En este curso trabajaremos con datos reales de estudiantes del programa de Ciencias Econ\'omicas y Administrativas}

\vspace{1em}
\begin{exampleblock}{Nuestros datos}
\begin{itemize}
\item Saber Pro 2023 y 2024 --- programa de Ciencias Econ\'omicas y Administrativas
\item Saber 11 enlazado para los mismos estudiantes
\item Almacenados en Google BigQuery (proyecto \texttt{ph-jabri})
\end{itemize}
\end{exampleblock}
\end{frame}

\begin{frame}{Variable central: MOD\_INGLES\_PUNT}
\begin{exampleblock}{¿Por qué el puntaje de inglés?}
Para estudiantes de Ciencias Econ\'omicas y Administrativas, la competencia en inglés es un \textbf{proxy de competitividad internacional}:
\begin{itemize}
\item Comunicación con clientes y proveedores internacionales
\item Acceso a información y literatura de negocios global
\item Movilidad laboral y oportunidades en empresas multinacionales
\item Negociación en mercados internacionales
\end{itemize}
\end{exampleblock}
\end{frame}

\begin{frame}{Preguntas que exploraremos}
\textbf{A lo largo del curso, usaremos \variable{MOD\_INGLES\_PUNT} para responder:}

\begin{enumerate}
\item ¿Cuál es el puntaje promedio en inglés de estudiantes de Ciencias Econ\'omicas?
\item ¿Cómo varía según tipo de institución (pública vs privada)?
\item ¿Cómo varía según estrato socioeconómico?
\item ¿Está relacionado con el desempeño en Saber 11?
\item ¿Hay diferencias regionales (por departamento)?
\end{enumerate}

\vspace{1em}
\alert{Hoy comenzamos con las preguntas 1--3, usando estadística descriptiva numérica.}
\end{frame}

\begin{frame}{Software: R y RStudio}
\begin{columns}[T]
\column{0.48\textwidth}
\textbf{¿Qué es R?}
\begin{itemize}
\item Lenguaje de programación estadística
\item Software libre y de código abierto
\item Más de 19,000 paquetes disponibles
\item Estándar en investigación y análisis de datos
\end{itemize}

\column{0.48\textwidth}
\textbf{¿Qué es RStudio?}
\begin{itemize}
\item Entorno de desarrollo integrado (IDE) para R
\item Editor de código, consola, gráficos, ayuda
\item Gestión de proyectos
\item Integración con Git, Markdown, Shiny
\end{itemize}
\end{columns}

\vspace{1em}
\textbf{Instalación:}
\begin{enumerate}
\item Descargar R desde \url{https://cran.r-project.org}
\item Descargar RStudio desde \url{https://posit.co/download/rstudio-desktop/}
\end{enumerate}
\end{frame}

\begin{frame}[fragile]{Primeros pasos en R}
\begin{lstlisting}
# Primeros comandos en R
2 + 3
x <- 10
x * 2

# Tipos basicos
nombre <- "Jaime"       # character
edad <- 25              # numeric
activo <- TRUE           # logical

# Vectores
puntajes <- c(145, 160, 155, 170, 150)
length(puntajes)
\end{lstlisting}
\end{frame}

\begin{frame}{Tipos de datos}
\begin{block}{Datos categóricos (cualitativos)}
Describen cualidades o características. No tienen un orden numérico natural.
\begin{itemize}
\item \textbf{Nominal}: sin orden (ej: tipo de IES [pública, privada], departamento)
\item \textbf{Ordinal}: con orden (ej: estrato [1, 2, 3, 4, 5, 6], nivel educativo)
\end{itemize}
\end{block}

\begin{block}{Datos cuantitativos (numéricos)}
Representan cantidades medibles.
\begin{itemize}
\item \textbf{Intervalo}: diferencias significativas, pero no hay cero absoluto (ej: temperatura en °C)
\item \textbf{Razón}: hay cero absoluto, se pueden hacer cocientes (ej: puntajes ICFES, edad, ingreso)
\end{itemize}
\end{block}
\end{frame}

\begin{frame}{Tipos de datos: ejemplos con ICFES}
\begin{alertblock}{Clasificación de variables del dataset}
\variable{MOD\_INGLES\_PUNT} es cuantitativa de razón (0 = ausencia total de conocimiento).\\
\variable{ESTU\_TIPOREMUNERACION} es categórica nominal.
\end{alertblock}

\vspace{1em}
\begin{table}
\centering
\small
\begin{tabular}{lll}
\toprule
\textbf{Variable} & \textbf{Tipo} & \textbf{Escala} \\
\midrule
\variable{MOD\_INGLES\_PUNT} & Cuantitativa & Razón \\
\variable{FAMI\_ESTRATOVIVIENDA} & Categórica & Ordinal \\
\variable{ESTU\_DEPTO\_RESIDE} & Categórica & Nominal \\
\variable{INST\_CARACTER\_ACADEMICO} & Categórica & Nominal \\
\bottomrule
\end{tabular}
\end{table}
\end{frame}

\begin{frame}{Datos de corte transversal vs series de tiempo}
\begin{columns}[T]
\column{0.48\textwidth}
\textbf{Corte transversal (cross-section)}
\begin{itemize}
\item Observaciones de múltiples unidades en un \alert{único momento en el tiempo}
\item Ejemplo: puntajes Saber Pro 2024 de todos los estudiantes de Ciencias Econ\'omicas
\item Nos permite comparar entre individuos, instituciones, regiones
\end{itemize}

\column{0.48\textwidth}
\textbf{Series de tiempo (time series)}
\begin{itemize}
\item Observaciones de una unidad a lo largo del \alert{tiempo}
\item Ejemplo: puntaje promedio en inglés de Ciencias Econ\'omicas entre 2015 y 2025
\item Nos permite identificar tendencias, estacionalidad
\end{itemize}
\end{columns}

\vspace{1em}
\begin{exampleblock}{En este curso}
Trabajaremos principalmente con \textbf{datos de corte transversal}: estudiantes de Ciencias Econ\'omicas en Saber Pro 2023 y 2024.
\end{exampleblock}
\end{frame}

% ============================================================================
\section{Medidas de tendencia central}
% ============================================================================

\begin{frame}{Media aritmética}
\begin{block}{Definición}
La \textbf{media aritmética} (o promedio) es la suma de los valores dividida por el número de observaciones.
\end{block}

\begin{columns}[T]
\column{0.48\textwidth}
\textbf{Media muestral} (estadístico):
\[
\bar{x} = \frac{1}{n} \sum_{i=1}^n x_i
\]
donde $n$ = tamaño de la muestra

\column{0.48\textwidth}
\textbf{Media poblacional} (parámetro):
\[
\mu = \frac{1}{N} \sum_{i=1}^N x_i
\]
donde $N$ = tamaño de la población
\end{columns}

\vspace{0.5em}
\alert{Notación}: $\bar{x}$ (x-barra) para muestra, $\mu$ (mu) para población

\begin{exampleblock}{Ejemplo: screen time semanal de 5 amigos}
Horas diarias: 5, 7, 4, 8, 6 $\;\Rightarrow\;$ $\bar{x} = (5+7+4+8+6)/5 = 6$ horas
\end{exampleblock}
\end{frame}

\begin{frame}{Media aritmética: ejemplo}
\begin{exampleblock}{Minutos de m\'usica en Spotify en 5 d\'ias}
Minutos escuchados: 45, 60, 55, 70, 50
\[
\bar{x} = \frac{45 + 60 + 55 + 70 + 50}{5} = \frac{280}{5} = 56 \text{ min}
\]
\end{exampleblock}

\vspace{1em}
\textbf{Propiedades de la media}:
\begin{itemize}
\item Usa toda la información disponible (todos los valores)
\item Es sensible a valores extremos (outliers)
\item La suma de desviaciones respecto a la media es cero: $\sum(x_i - \bar{x}) = 0$
\end{itemize}
\end{frame}

\begin{frame}[fragile]{Media aritmética en R}
\begin{lstlisting}
# Vector de puntajes
puntajes <- c(145, 160, 155, 170, 150, 162, 158, 149, 165, 153)

# Calcular la media
mean(puntajes)
# [1] 156.7

# Con datos faltantes (NA), usar na.rm = TRUE
puntajes_na <- c(145, 160, NA, 170, 150)
mean(puntajes_na, na.rm = TRUE)
# [1] 156.25

# Cargar datos ICFES completos
library(tidyverse)
icfes <- read_csv("datos/saber_pro_ni_2024.csv")
mean(icfes$MOD_INGLES_PUNT, na.rm = TRUE)
\end{lstlisting}

\alert{Nota}: \code{na.rm = TRUE} elimina valores faltantes antes de calcular
\end{frame}

\begin{frame}{Media ponderada}
\begin{block}{Definición}
La \textbf{media ponderada} asigna pesos diferentes a cada observación según su importancia.
\end{block}

\[
\bar{x}_w = \frac{\sum_{i=1}^n w_i x_i}{\sum_{i=1}^n w_i}
\]

donde $w_i$ es el peso de la observación $i$
\end{frame}

\begin{frame}{Media ponderada: ejemplo}
\begin{exampleblock}{Tu Spotify Wrapped: top 3 artistas}
Tus artistas m\'as escuchados con distinto \% de tu tiempo total:
\begin{itemize}
\item Bad Bunny: rating 9.2 (50\% de tu tiempo)
\item Karol G: rating 8.5 (30\% de tu tiempo)
\item Feid: rating 7.8 (20\% de tu tiempo)
\end{itemize}
\[
\bar{x}_w = \frac{0.50(9.2) + 0.30(8.5) + 0.20(7.8)}{1} = 8.71
\]
\end{exampleblock}

\vspace{0.5em}
\alert{La media simple dar\'ia}: $(9.2 + 8.5 + 7.8)/3 = 8.50$. La ponderada refleja que escuchas m\'as a Bad Bunny.
\end{frame}

\begin{frame}{Mediana}
\begin{block}{Definición}
La \textbf{mediana} es el valor que divide los datos ordenados en dos mitades iguales. 50\% de las observaciones están por debajo y 50\% por encima.
\end{block}

\textbf{Procedimiento:}
\begin{enumerate}
\item Ordenar los datos de menor a mayor
\item Si $n$ es impar: la mediana es el valor central
\item Si $n$ es par: la mediana es el promedio de los dos valores centrales
\end{enumerate}

\begin{columns}[T]
\column{0.48\textwidth}
\textbf{Ejemplo (n impar):}\\
\footnotesize
Seguidores en IG de 5 amigos:\\
280, 310, \alert{350}, 400, 520\\
Mediana = 350

\column{0.48\textwidth}
\textbf{Ejemplo (n par):}\\
\footnotesize
Seguidores de 6 amigos:\\
280, 310, \alert{350, 400}, 520, 610\\
Mediana = (350+400)/2 = 375
\end{columns}

\begin{alertblock}{Propiedad clave}
\small La mediana es \textbf{robusta a outliers}. Si un amigo tiene 50{,}000 seguidores porque se hizo viral, la media sube a locura, pero la mediana ni se entera.
\end{alertblock}
\end{frame}

\begin{frame}[fragile]{Mediana en R}
\begin{lstlisting}
# Vector de puntajes
puntajes <- c(145, 150, 155, 160, 170)
median(puntajes)
# [1] 155

# Puntajes con un valor extremo
puntajes_outlier <- c(145, 150, 155, 160, 250)
mean(puntajes_outlier)    # [1] 172 (afectada por el outlier)
median(puntajes_outlier)  # [1] 155 (no afectada)

# Con datos ICFES
median(icfes$MOD_INGLES_PUNT, na.rm = TRUE)

# Comparar media vs mediana
summary(icfes$MOD_INGLES_PUNT)
\end{lstlisting}

\vspace{0.5em}
\alert{Interpretación}: Si media > mediana, distribución sesgada a la derecha (cola larga hacia valores altos). Si media < mediana, sesgada a la izquierda.
\end{frame}

\begin{frame}{Moda}
\begin{block}{Definición}
La \textbf{moda} es el valor que aparece con mayor frecuencia en los datos.
\end{block}

\begin{itemize}
\item Datos pueden ser \textbf{unimodales} (una moda), \textbf{bimodales} (dos modas), o \textbf{multimodales}
\item Es la \'unica medida de tendencia central aplicable a datos categ\'oricos
\item La canci\'on m\'as escuchada de tu playlist este mes = la moda
\end{itemize}

\begin{columns}[T]
\column{0.48\textwidth}
\begin{exampleblock}{Unimodal}
150, 155, 155, 160, 155, 165, 170, 155\\
Moda = 155 (aparece 4 veces)
\end{exampleblock}

\column{0.48\textwidth}
\begin{exampleblock}{Bimodal}
2, 2, 2, 3, 4, 5, 5, 5, 6\\
Modas = 2 y 5 (3 veces c/u)
\end{exampleblock}
\end{columns}
\end{frame}

\begin{frame}[fragile]{Moda en R}
\begin{lstlisting}[basicstyle=\ttfamily\scriptsize]
# R no tiene una funcion integrada para moda
# Opcion 1: usar table() para ver frecuencias
puntajes <- c(150, 155, 155, 160, 155, 165, 170, 155)
table(puntajes)
# puntajes
# 150 155 160 165 170
#   1   4   1   1   1

# Opcion 2: funcion personalizada
moda <- function(x) {
  freq <- table(x)
  as.numeric(names(freq)[which.max(freq)])
}
moda(puntajes)
# [1] 155

# Para categoricas
table(icfes$ESTU_DEPTO_RESIDE)
# Muestra frecuencia por departamento
\end{lstlisting}
\end{frame}

\begin{frame}{¿Cuándo usar cada medida?}
\begin{table}
\centering
\small
\begin{tabular}{lp{6cm}p{5cm}}
\toprule
\textbf{Medida} & \textbf{Cuándo usar} & \textbf{Ventajas / Desventajas} \\
\midrule
Media & Distribuciones simétricas, datos numéricos & Usa toda la información; sensible a outliers \\
\addlinespace
Mediana & Distribuciones sesgadas, presencia de outliers & Robusta; ignora valores extremos \\
\addlinespace
Moda & Datos categóricos, distribuciones multimodales & Única para categorías; puede no existir o no ser única \\
\bottomrule
\end{tabular}
\end{table}

\vspace{0.3em}
\begin{alertblock}{Regla pr\'actica}
\begin{itemize}
\item Distribuci\'on sim\'etrica: media $\approx$ mediana $\approx$ moda
\item Sesgada a la derecha: moda < mediana < media (ej: likes en TikTok --- muchos posts con pocos likes, pocos virales)
\item Sesgada a la izquierda: media < mediana < moda
\end{itemize}
\end{alertblock}
\end{frame}

% ============================================================================
\section{Medidas de dispersión}
% ============================================================================

\begin{frame}{Rango}
\begin{block}{Definición}
El \textbf{rango} es la diferencia entre el valor máximo y el valor mínimo.
\[
\text{Rango} = \max(x) - \min(x)
\]
\end{block}

\begin{itemize}
\item Medida más simple de dispersión
\item Muy sensible a outliers (basta un valor extremo para inflarlo)
\item No usa información de los valores intermedios
\end{itemize}

\vspace{0.5em}
\begin{columns}[T]
\column{0.48\textwidth}
\begin{exampleblock}{Screen time del grupo}
Horas: 2, 4, 5, 7, 9\\
Rango = $9 - 2 = 7$\,h
\end{exampleblock}

\column{0.48\textwidth}
\begin{exampleblock}{Un amigo dej\'o el cel en casa}
Horas: \alert{0}, 4, 5, 7, 9\\
Rango = $9 - 0 = 9$\,h
\end{exampleblock}
\end{columns}
\end{frame}

\begin{frame}{Rango intercuart\'ilico (IQR)}
\begin{block}{Definici\'on}
El \textbf{rango intercuart\'ilico} (IQR) es la diferencia entre el tercer cuartil (Q3) y el primer cuartil (Q1). Contiene el 50\% central de los datos.
\[
\text{IQR} = Q_3 - Q_1
\]
\end{block}

\begin{itemize}
\item Robusto a outliers (ignora el 25\% m\'as bajo y el 25\% m\'as alto)
\item Usado en box plots para identificar outliers
\item Outlier si: $x < Q_1 - 1.5 \times \text{IQR}$ o $x > Q_3 + 1.5 \times \text{IQR}$
\end{itemize}
\end{frame}

\begin{frame}{Rango intercuart\'ilico (IQR) (cont.)}
\begin{exampleblock}{Ejemplo: screen time diario (horas) de tu grupo}
Si Q1 = 3\,h y Q3 = 7\,h:\\
IQR = 7 - 3 = 4\,h\\
L\'imite inferior outlier: $3 - 1.5(4) = -3$ (no aplica)\\
L\'imite superior outlier: $7 + 1.5(4) = 13$\,h $\;\Rightarrow\;$ $>$13\,h ser\'ia at\'ipico
\end{exampleblock}
\end{frame}

\begin{frame}[fragile]{Rango e IQR en R}
\begin{lstlisting}[basicstyle=\ttfamily\scriptsize]
# Rango
puntajes <- c(145, 150, 155, 160, 165, 170)
range(puntajes)
# [1] 145 170
diff(range(puntajes))
# [1] 25

# Rango intercuartilico
IQR(puntajes)
# [1] 15

# Cuartiles
quantile(puntajes, probs = c(0.25, 0.75))
#   25%   75%
# 151.25 163.75

# Con datos ICFES
summary(icfes$MOD_INGLES_PUNT)
IQR(icfes$MOD_INGLES_PUNT, na.rm = TRUE)
\end{lstlisting}
\end{frame}

\begin{frame}{Varianza}
\begin{block}{Definición}
La \textbf{varianza} mide la dispersión promedio de los datos respecto a la media. Es el promedio de las desviaciones cuadradas.
\end{block}

\begin{columns}[T]
\column{0.48\textwidth}
\textbf{Varianza muestral}:
\[
s^2 = \frac{1}{n-1} \sum_{i=1}^n (x_i - \bar{x})^2
\]

\column{0.48\textwidth}
\textbf{Varianza poblacional}:
\[
\sigma^2 = \frac{1}{N} \sum_{i=1}^N (x_i - \mu)^2
\]
\end{columns}

\vspace{0.5em}
\textbf{Propiedades}: siempre $s^2 \geq 0$; unidades en cuadrado de las originales.

\begin{exampleblock}{Intuici\'on}
Un artista de Spotify con streams diarios estables (poca varianza) vs uno que tuvo un hit viral un d\'ia y nada al siguiente (mucha varianza).
\end{exampleblock}
\end{frame}

\begin{frame}{Varianza: grados de libertad}
\begin{alertblock}{¿Por qué $n-1$ y no $n$?}
Dividir por $n-1$ (grados de libertad) hace que $s^2$ sea un \textbf{estimador insesgado} de $\sigma^2$. Dividir por $n$ subestimaría la varianza poblacional sistemáticamente.
\end{alertblock}

\vspace{1em}
\textbf{Intuición}:
\begin{itemize}
\item Al estimar $\sigma^2$, usamos $\bar{x}$ en lugar de $\mu$
\item Esto ``consume'' un grado de libertad
\item Quedan $n-1$ piezas de información independientes
\item La corrección compensa el sesgo de usar una estimación ($\bar{x}$) en lugar del valor real ($\mu$)
\end{itemize}
\end{frame}

\begin{frame}{Desviación estándar}
\begin{block}{Definición}
La \textbf{desviación estándar} es la raíz cuadrada positiva de la varianza. Mide la dispersión en las mismas unidades que los datos.
\end{block}

\begin{columns}[T]
\column{0.48\textwidth}
\textbf{Desviación estándar muestral}:
\[
s = \sqrt{s^2} = \sqrt{\frac{1}{n-1} \sum_{i=1}^n (x_i - \bar{x})^2}
\]

\column{0.48\textwidth}
\textbf{Desviación estándar poblacional}:
\[
\sigma = \sqrt{\sigma^2} = \sqrt{\frac{1}{N} \sum_{i=1}^N (x_i - \mu)^2}
\]
\end{columns}

\vspace{1em}
\textbf{Interpretación}:
\begin{itemize}
\item Mide cu\'anto se desv\'ian los datos de la media \alert{en promedio}
\item $s$ peque\~na: datos concentrados cerca de $\bar{x}$ (screen time estable cada d\'ia)
\item $s$ grande: datos dispersos, alejados de $\bar{x}$ (un d\'ia 1\,h, otro d\'ia 10\,h)
\item Mismas unidades que los datos originales (m\'as interpretable que $s^2$)
\end{itemize}
\end{frame}

\begin{frame}[fragile]{Varianza y desviación estándar en R (1/2)}
\begin{lstlisting}
puntajes <- c(145, 150, 155, 160, 165, 170)

# Varianza (usa n-1)
var(puntajes)
# [1] 104.1667

# Desviacion estandar
sd(puntajes)
# [1] 10.20621

# Verificar que sd = sqrt(var)
sqrt(var(puntajes))
# [1] 10.20621
\end{lstlisting}
\end{frame}

\begin{frame}[fragile]{Varianza y desviación estándar en R (2/2)}
\begin{lstlisting}
# Calculo manual (para entender)
puntajes <- c(145, 150, 155, 160, 165, 170)
n <- length(puntajes)
media <- mean(puntajes)
varianza_manual <- sum((puntajes - media)^2) / (n - 1)
varianza_manual
# [1] 104.1667

# Con datos ICFES
sd(icfes$MOD_INGLES_PUNT, na.rm = TRUE)
\end{lstlisting}
\end{frame}

\begin{frame}{Coeficiente de variación}
\begin{block}{Definición}
El \textbf{coeficiente de variación} (CV) expresa la desviación estándar como porcentaje de la media. Compara dispersión relativa entre variables con distintas escalas.
\[
CV = \frac{s}{\bar{x}} \times 100\%
\]
\end{block}

\begin{columns}[T]
\column{0.55\textwidth}
\textbf{Interpretación}:
\begin{itemize}
\item CV $< 15\%$: datos homogéneos
\item CV $15$--$30\%$: variabilidad moderada
\item CV $> 30\%$: datos heterogéneos
\end{itemize}

\column{0.4\textwidth}
\begin{exampleblock}{Ejemplo}
$\bar{x} = 157$, $s = 20$\\[0.3em]
$CV = \frac{20}{157} \times 100\% = 12.7\%$
\end{exampleblock}
\end{columns}

\vspace{0.3em}
\alert{Sólo válido para variables de razón} ($\bar{x} > 0$)
\end{frame}

\begin{frame}[fragile]{Coeficiente de variación en R}
\begin{lstlisting}
# R no tiene funcion integrada para CV
# Funcion personalizada
cv <- function(x, na.rm = FALSE) {
  (sd(x, na.rm = na.rm) / mean(x, na.rm = na.rm)) * 100
}

puntajes <- c(145, 150, 155, 160, 165, 170)
cv(puntajes)
# [1] 6.52

# Comparar CV de ingles vs matematicas
cv_ingles <- cv(icfes$MOD_INGLES_PUNT, na.rm = TRUE)
cv_matematicas <- cv(icfes$MOD_RAZONA_CUANTITAT_PUNT, na.rm = TRUE)

cat("CV Ingles:", round(cv_ingles, 2), "%\n")
cat("CV Matematicas:", round(cv_matematicas, 2), "%\n")

# ?Cual tiene mayor variabilidad relativa?
\end{lstlisting}
\end{frame}

% ============================================================================
\section{Forma de la distribución}
% ============================================================================

\begin{frame}{Puntajes estandarizados (z-scores)}
\begin{block}{Definición}
El \textbf{puntaje estandarizado} o \textbf{z-score} indica cuántas desviaciones estándar se encuentra una observación por encima (o por debajo) de la media.
\[
z_i = \frac{x_i - \bar{x}}{s}
\]
\end{block}

\textbf{Propiedades}:
\begin{itemize}
\item Sin unidades (adimensional)
\item Media de los z-scores = 0, desviación estándar = 1
\item $z > 0$: por encima de la media; $z < 0$: por debajo
\item $|z| > 2$: observación inusual; $|z| > 3$: posible outlier
\end{itemize}
\end{frame}

\begin{frame}{Z-scores: ejemplo}
\begin{exampleblock}{Ejemplo: screen time diario entre tus amigos}
Si $\bar{x} = 5$\,h, $s = 1.2$\,h, y tu screen time es $x = 8$\,h:
\[
z = \frac{8 - 5}{1.2} = 2.5
\]
Est\'as 2.5 desviaciones est\'andar por encima de tus amigos. Eres un outlier de celular.
\end{exampleblock}

\vspace{1em}
\begin{alertblock}{Utilidad}
Los z-scores permiten:
\begin{itemize}
\item Comparar observaciones de variables con distintas escalas
\item Identificar valores atípicos de forma objetiva
\item Estandarizar datos antes de análisis multivariados
\end{itemize}
\end{alertblock}
\end{frame}

\begin{frame}[fragile]{Z-scores en R (1/2)}
\begin{lstlisting}
puntajes <- c(130, 140, 150, 155, 160, 170, 185)

# Funcion para calcular z-scores
z_scores <- function(x) {
  (x - mean(x)) / sd(x)
}

z <- z_scores(puntajes)
z
# [1] -1.456  -0.873  -0.291  0.000  0.291  0.873  1.456

# Verificar propiedades
mean(z)  # [1] 0 (aproximadamente)
sd(z)    # [1] 1
\end{lstlisting}
\end{frame}

\begin{frame}[fragile]{Z-scores en R (2/2)}
\begin{lstlisting}
# Identificar outliers (|z| > 3)
puntajes <- c(130, 140, 150, 155, 160, 170, 185)
z <- (puntajes - mean(puntajes)) / sd(puntajes)
outliers <- puntajes[abs(z) > 3]
outliers

# Con datos ICFES
icfes <- icfes %>%
  mutate(z_ingles = (MOD_INGLES_PUNT - mean(MOD_INGLES_PUNT, na.rm=T)) /
                     sd(MOD_INGLES_PUNT, na.rm=T))
\end{lstlisting}
\end{frame}

\begin{frame}{Teorema de Chebyshev}
\begin{block}{Teorema de Chebyshev}
Para \alert{cualquier distribución}, al menos $1 - \frac{1}{k^2}$ de las observaciones caen dentro de $k$ desviaciones estándar de la media, donde $k > 1$.
\end{block}

\begin{table}
\centering
\begin{tabular}{ccc}
\toprule
$k$ & Intervalo & \% mínimo de datos \\
\midrule
1.5 & $\bar{x} \pm 1.5s$ & $1 - 1/(1.5)^2 = 55.6\%$ \\
2 & $\bar{x} \pm 2s$ & $1 - 1/4 = 75\%$ \\
3 & $\bar{x} \pm 3s$ & $1 - 1/9 = 88.9\%$ \\
4 & $\bar{x} \pm 4s$ & $1 - 1/16 = 93.75\%$ \\
\bottomrule
\end{tabular}
\end{table}

\vspace{0.5em}
\begin{alertblock}{Importante}
Este teorema aplica a \textbf{cualquier distribución} (no requiere normalidad). Es una cota conservadora: distribuciones específicas pueden tener proporciones mayores.
\end{alertblock}
\end{frame}

\begin{frame}{Regla empírica (68-95-99.7)}
\begin{block}{Regla empírica}
Para distribuciones \alert{aproximadamente normales} (en forma de campana):
\begin{itemize}
\item Aproximadamente 68\% de los datos caen dentro de $\bar{x} \pm 1s$
\item Aproximadamente 95\% de los datos caen dentro de $\bar{x} \pm 2s$
\item Aproximadamente 99.7\% de los datos caen dentro de $\bar{x} \pm 3s$
\end{itemize}
\end{block}

\begin{center}
\begin{tikzpicture}[scale=0.8]
\begin{axis}[
  domain=-4:4, samples=100,
  axis lines=middle, xlabel={$z$}, ylabel={},
  ytick=\empty, xtick={-3,-2,-1,0,1,2,3},
  xticklabels={$-3s$,$-2s$,$-1s$,$\bar{x}$,$1s$,$2s$,$3s$},
  height=4cm, width=11cm,
  every axis x label/.style={at={(ticklabel* cs:1)},anchor=west},
  enlargelimits=false, clip=false
]
\addplot[very thick, javeriana] {exp(-x^2/2)/sqrt(2*pi)};
\addplot[fill=javerianalight, opacity=0.3, domain=-1:1] {exp(-x^2/2)/sqrt(2*pi)} \closedcycle;
\node at (0, 0.2) {68\%};
\addplot[fill=javerianalight, opacity=0.15, domain=-2:-1] {exp(-x^2/2)/sqrt(2*pi)} \closedcycle;
\addplot[fill=javerianalight, opacity=0.15, domain=1:2] {exp(-x^2/2)/sqrt(2*pi)} \closedcycle;
\node at (1.5, 0.05) {\tiny 13.5\%};
\node at (-1.5, 0.05) {\tiny 13.5\%};
\end{axis}
\end{tikzpicture}
\end{center}

\alert{Importante}: Sólo aplica a distribuciones aproximadamente normales.
\end{frame}

\begin{frame}{Asimetr\'ia (skewness)}
\begin{block}{Definici\'on}
La \textbf{asimetr\'ia} mide el grado de simetr\'ia de una distribuci\'on.
\[
g_1 = \frac{1}{n} \sum_{i=1}^n \left(\frac{x_i - \bar{x}}{s}\right)^3
\]
\end{block}

\begin{columns}[T]
\column{0.33\textwidth}
\textbf{Sim\'etrica} ($g_1 \approx 0$)\\
\textit{Visualizaci\'on generada en R con} \code{plot()} \textit{y} \code{hist()}\\
Media = Mediana

\column{0.33\textwidth}
\textbf{Asimetr\'ia positiva} ($g_1 > 0$)\\
Cola larga a la derecha\\
Media > Mediana\\
\textit{Ej: likes en posts --- la mayor\'ia recibe pocos}

\column{0.33\textwidth}
\textbf{Asimetr\'ia negativa} ($g_1 < 0$)\\
Cola larga a la izquierda\\
Media < Mediana
\end{columns}
\end{frame}

\begin{frame}{Asimetr\'ia (skewness) (cont.)}
\alert{Interpretaci\'on}:
\begin{itemize}
\item $|g_1| < 0.5$: aproximadamente sim\'etrica
\item $0.5 \leq |g_1| < 1$: asimetr\'ia moderada
\item $|g_1| \geq 1$: asimetr\'ia pronunciada
\end{itemize}
\end{frame}

\begin{frame}{Curtosis}
\begin{block}{Definición}
La \textbf{curtosis} mide el peso de las colas y la altura del pico de una distribución, comparado con la distribución normal.
\[
g_2 = \frac{1}{n} \sum_{i=1}^n \left(\frac{x_i - \bar{x}}{s}\right)^4 - 3
\]
(El $-3$ es para que la distribución normal tenga curtosis = 0)
\end{block}

\begin{itemize}
\item \textbf{Leptocúrtica} ($g_2 > 0$): colas pesadas, pico alto, más concentrada que la normal
\item \textbf{Mesocúrtica} ($g_2 \approx 0$): similar a la normal
\item \textbf{Platocúrtica} ($g_2 < 0$): colas ligeras, pico bajo, más plana que la normal
\end{itemize}

\vspace{0.5em}
\textbf{Importancia práctica}:
\begin{itemize}
\item $g_2 > 0$ indica mayor presencia de valores extremos (riesgo financiero)
\item Afecta la validez de pruebas estadísticas que asumen normalidad
\end{itemize}
\end{frame}

\begin{frame}[fragile]{Asimetría y curtosis en R (1/2)}
\begin{lstlisting}
# Instalar y cargar paquete moments
install.packages("moments")
library(moments)

puntajes <- c(130, 135, 140, 145, 150, 155, 160, 165, 190)

# Asimetria
skewness(puntajes)
# [1] 0.847 (asimetria positiva moderada)

# Curtosis (exceso de curtosis)
kurtosis(puntajes) - 3
# [1] -0.312 (aproximadamente mesocurtica)
\end{lstlisting}
\end{frame}

\begin{frame}[fragile]{Asimetría y curtosis en R (2/2)}
\begin{lstlisting}
# Con datos ICFES
skewness(icfes$MOD_INGLES_PUNT, na.rm = TRUE)
kurtosis(icfes$MOD_INGLES_PUNT, na.rm = TRUE) - 3

# Interpretacion
if (skewness(icfes$MOD_INGLES_PUNT, na.rm=T) > 0) {
  print("Distribucion sesgada a la derecha")
}
\end{lstlisting}
\end{frame}

\begin{frame}{Detección de outliers: métodos}
\textbf{1. Método IQR (box plot)}:
\begin{itemize}
\item Outlier si: $x < Q_1 - 1.5 \times \text{IQR}$ o $x > Q_3 + 1.5 \times \text{IQR}$
\item Outlier extremo si: $x < Q_1 - 3 \times \text{IQR}$ o $x > Q_3 + 3 \times \text{IQR}$
\end{itemize}

\textbf{2. Método z-score}:
\begin{itemize}
\item Inusual si: $|z| > 2$
\item Outlier si: $|z| > 3$
\end{itemize}

\textbf{3. Método de desviación absoluta mediana (MAD)}:
\begin{itemize}
\item Más robusto que z-score para distribuciones no normales
\end{itemize}
\end{frame}

\begin{frame}{¿Qué hacer con outliers?}
\begin{alertblock}{Protocolo para tratar valores atípicos}
\begin{enumerate}
\item Verificar si es un error de medición o entrada de datos
\item Investigar si representa un fenómeno real
\item Considerar análisis con y sin outliers
\item Nunca eliminar sin justificación
\end{enumerate}
\end{alertblock}

\vspace{1em}
\begin{exampleblock}{Ejemplo}
Ese post que se hizo viral con 50K likes cuando normalmente recibes 200. ?`Es un error? No, es un fen\'omeno real. Pero si tu dataset dice que un post tuvo $-30$ likes, eso s\'i es un error de datos.
\end{exampleblock}
\end{frame}

% ============================================================================
\section{Percentiles y cuartiles}
% ============================================================================

\begin{frame}{Percentiles}
\begin{block}{Definición}
El \textbf{percentil $p$} es el valor tal que al menos el $p\%$ de las observaciones son menores o iguales a él.
\end{block}

\textbf{Percentiles clave}:
\begin{itemize}
\item Percentil 25 (P$_{25}$) = Q$_1$: 25\% de los datos están por debajo
\item Percentil 50 (P$_{50}$) = Q$_2$ = mediana
\item Percentil 75 (P$_{75}$) = Q$_3$: 75\% de los datos están por debajo
\item Percentil 90 (P$_{90}$): 90\% de los datos están por debajo
\end{itemize}

\vspace{0.5em}
\textbf{Deciles}: dividen los datos en 10 partes (D$_1$, D$_2$, \ldots, D$_9$)\\
\textbf{Quintiles}: dividen los datos en 5 partes
\end{frame}

\begin{frame}{Percentiles: interpretación}
\begin{exampleblock}{Ejemplo: Spotify Wrapped}
Spotify Wrapped te dice que est\'as en el \textbf{top 1\% de oyentes de Bad Bunny}. Eso es el percentil 99: escuchaste m\'as que el 99\% de sus oyentes.
\end{exampleblock}

\vspace{1em}
\begin{alertblock}{Nota}
Los percentiles son especialmente útiles para:
\begin{itemize}
\item Comparar el rendimiento relativo de un estudiante
\item Establecer puntos de corte (ej: becas para el top 10\%)
\item Describir la distribución sin asumir normalidad
\end{itemize}
\end{alertblock}
\end{frame}

\begin{frame}{Cuartiles y five-number summary}
\begin{block}{Cuartiles}
Los \textbf{cuartiles} dividen los datos ordenados en cuatro partes iguales:
\begin{itemize}
\item Q$_1$: primer cuartil (25\%)
\item Q$_2$: segundo cuartil (50\%) = mediana
\item Q$_3$: tercer cuartil (75\%)
\end{itemize}
\end{block}

\begin{block}{Five-number summary}
Resumen de 5 números que describe la distribución:
\[
\{\text{mín}, Q_1, Q_2, Q_3, \text{máx}\}
\]
\end{block}
\end{frame}

\begin{frame}{Five-number summary: ejemplo}
\begin{exampleblock}{Puntajes de inglés (Saber Pro)}
Five-number summary: $\{100, 145, 157, 168, 200\}$
\begin{itemize}
\item 25\% de estudiantes obtuvo $\leq$ 145 puntos
\item 50\% obtuvo $\leq$ 157 puntos (mediana)
\item 75\% obtuvo $\leq$ 168 puntos
\end{itemize}
\end{exampleblock}

\vspace{1em}
\alert{El five-number summary es la base del box plot}, que veremos en la próxima sesión.
\end{frame}

\begin{frame}[fragile]{Percentiles y cuartiles en R (1/2)}
\begin{lstlisting}
puntajes <- c(120, 130, 140, 145, 150, 155, 160, 165, 170, 180)

# Cuartiles
quantile(puntajes)
#   0%  25%  50%  75% 100%
#  120  141  152  163  180

# Percentiles especificos
quantile(puntajes, probs = c(0.10, 0.25, 0.50, 0.75, 0.90))
#  10%  25%  50%  75%  90%
#  130  141  152  163  174
\end{lstlisting}
\end{frame}

\begin{frame}[fragile]{Percentiles y cuartiles en R (2/2)}
\begin{lstlisting}
# Five-number summary
fivenum(puntajes)
# [1] 120 140 152 165 180

# Summary completo (incluye media)
summary(puntajes)
#  Min. 1st Qu.  Median    Mean 3rd Qu.    Max.
#   120     141     152     151     163     180

# Con datos ICFES
summary(icfes$MOD_INGLES_PUNT)
\end{lstlisting}
\end{frame}

% ============================================================================
\section{Aplicación con datos ICFES}
% ============================================================================

\begin{frame}[fragile]{Ejemplo completo: cargar y explorar datos}
\begin{lstlisting}[basicstyle=\ttfamily\scriptsize]
# Cargar librerias
library(tidyverse)

# Cargar datos
icfes <- read_csv("datos/saber_pro_ni_2024.csv")

# Dimension de los datos
dim(icfes)
# [1] 3245   42

# Primeras filas
head(icfes)

# Estructura
str(icfes)

# Resumen de variables numericas
summary(icfes %>% select(MOD_INGLES_PUNT, MOD_LECTURA_CRITICA_PUNT,
                          MOD_RAZONA_CUANTITAT_PUNT))
\end{lstlisting}
\end{frame}

\begin{frame}[fragile]{Descriptivas de MOD\_INGLES\_PUNT: tendencia central}
\begin{lstlisting}
# Medidas de tendencia central
media_ingles <- mean(icfes$MOD_INGLES_PUNT, na.rm = TRUE)
mediana_ingles <- median(icfes$MOD_INGLES_PUNT, na.rm = TRUE)

cat("Media:", round(media_ingles, 2), "\n")
cat("Mediana:", round(mediana_ingles, 2), "\n")

# Medidas de dispersion
sd_ingles <- sd(icfes$MOD_INGLES_PUNT, na.rm = TRUE)
var_ingles <- var(icfes$MOD_INGLES_PUNT, na.rm = TRUE)
cv_ingles <- (sd_ingles / media_ingles) * 100

cat("Desviacion estandar:", round(sd_ingles, 2), "\n")
cat("Varianza:", round(var_ingles, 2), "\n")
cat("CV:", round(cv_ingles, 2), "%\n")
\end{lstlisting}
\end{frame}

\begin{frame}[fragile]{Descriptivas de MOD\_INGLES\_PUNT: forma}
\begin{lstlisting}
# Forma de la distribucion
library(moments)
skew_ingles <- skewness(icfes$MOD_INGLES_PUNT, na.rm = TRUE)
kurt_ingles <- kurtosis(icfes$MOD_INGLES_PUNT, na.rm = TRUE) - 3

cat("Asimetria:", round(skew_ingles, 2), "\n")
cat("Curtosis:", round(kurt_ingles, 2), "\n")
\end{lstlisting}

\vspace{0.5em}
\alert{Interpretación}: Si la asimetría es positiva, hay una cola larga hacia puntajes altos. Si la curtosis es positiva, hay más valores extremos de lo esperado bajo normalidad.
\end{frame}

\begin{frame}[fragile]{Desagregaci\'on por tipo de IES}
\begin{lstlisting}[basicstyle=\ttfamily\scriptsize]
# Comparar puntajes por tipo de institucion
icfes %>%
  group_by(INST_CARACTER_ACADEMICO) %>%
  summarise(
    n = n(),
    media = mean(MOD_INGLES_PUNT, na.rm = TRUE),
    mediana = median(MOD_INGLES_PUNT, na.rm = TRUE),
    sd = sd(MOD_INGLES_PUNT, na.rm = TRUE),
    cv = (sd / media) * 100,
    min = min(MOD_INGLES_PUNT, na.rm = TRUE),
    max = max(MOD_INGLES_PUNT, na.rm = TRUE)
  ) %>%
  arrange(desc(media))

# Resultado esperado:
# INST_CARACTER_ACADEMICO    n  media mediana    sd    cv  min  max
# Universidad (privada)    1850  162.3   161.0  18.5  11.4  110  200
# Universidad (publica)     950  157.8   156.0  19.2  12.2  105  198
# Inst. Tecnica/Tecnologica 445  148.5   147.0  20.1  13.5  100  195
\end{lstlisting}
\end{frame}

\begin{frame}[fragile]{Desagregación por estrato}
\begin{lstlisting}
# Comparar puntajes por estrato socioeconomico
icfes %>%
  filter(!is.na(FAMI_ESTRATOVIVIENDA)) %>%
  group_by(FAMI_ESTRATOVIVIENDA) %>%
  summarise(
    n = n(),
    media = mean(MOD_INGLES_PUNT, na.rm = TRUE),
    mediana = median(MOD_INGLES_PUNT, na.rm = TRUE),
    sd = sd(MOD_INGLES_PUNT, na.rm = TRUE),
    Q1 = quantile(MOD_INGLES_PUNT, 0.25, na.rm = TRUE),
    Q3 = quantile(MOD_INGLES_PUNT, 0.75, na.rm = TRUE),
    IQR = Q3 - Q1
  ) %>%
  arrange(FAMI_ESTRATOVIVIENDA)

# ?Hay una relacion clara entre estrato y puntaje de ingles?
# ?Como varia la dispersion entre estratos?
\end{lstlisting}
\end{frame}

\begin{frame}[fragile]{Tabla resumen: función personalizada}
\begin{lstlisting}
# Funcion para calcular todas las estadisticas
estadisticas_completas <- function(x) {
  tibble(
    n = sum(!is.na(x)),
    media = mean(x, na.rm = TRUE),
    mediana = median(x, na.rm = TRUE),
    sd = sd(x, na.rm = TRUE),
    cv = (sd / media) * 100,
    min = min(x, na.rm = TRUE),
    Q1 = quantile(x, 0.25, na.rm = TRUE),
    Q3 = quantile(x, 0.75, na.rm = TRUE),
    max = max(x, na.rm = TRUE),
    IQR = Q3 - Q1,
    asimetria = skewness(x, na.rm = TRUE),
    curtosis = kurtosis(x, na.rm = TRUE) - 3
  )
}
\end{lstlisting}
\end{frame}

\begin{frame}[fragile]{Tabla resumen: aplicación}
\begin{lstlisting}
# Aplicar a puntaje de ingles
estadisticas_completas(icfes$MOD_INGLES_PUNT)

# Aplicar a multiples variables
variables <- icfes %>%
  select(MOD_INGLES_PUNT, MOD_LECTURA_CRITICA_PUNT,
         MOD_RAZONA_CUANTITAT_PUNT)

variables %>%
  summarise(across(everything(), ~estadisticas_completas(.x)))
\end{lstlisting}

\vspace{0.5em}
\alert{Tip}: Esta función reutilizable les servirá en los Problem Sets y el proyecto final.
\end{frame}

% ============================================================================
\section{Resumen y próxima sesión}
% ============================================================================

\begin{frame}{Resumen de la sesión}
\textbf{Conceptos clave cubiertos}:

\begin{enumerate}
\item \textbf{Introducción}: estadística descriptiva vs inferencial, datos ICFES, tipos de datos
\item \textbf{Tendencia central}: media, mediana, moda
\item \textbf{Dispersión}: rango, IQR, varianza, desviación estándar, CV
\item \textbf{Forma}: z-scores, Chebyshev, regla empírica, asimetría, curtosis
\item \textbf{Posición}: percentiles, cuartiles, five-number summary
\item \textbf{Aplicación}: análisis completo de MOD\_INGLES\_PUNT
\end{enumerate}

\vspace{1em}
\begin{alertblock}{Para la próxima sesión}
\begin{itemize}
\item Repasar funciones de R: \code{mean()}, \code{median()}, \code{sd()}, \code{quantile()}, \code{summary()}
\item Traer computadora con R y RStudio instalados
\item Descargar los datos ICFES del repositorio del curso
\end{itemize}
\end{alertblock}
\end{frame}

\begin{frame}[fragile]{Próxima sesión: Estadística Descriptiva Gráfica}
\textbf{Vie 14 ago 2026}

\textbf{Temas a cubrir}:
\begin{itemize}
\item Distribuciones de frecuencia
\item Visualización de una variable: histogramas, box plots, density plots
\item Visualización de dos variables: scatter plots, heatmaps
\item Covarianza y correlación (Pearson y Spearman)
\item Paradoja de Simpson
\item Pipeline completo de análisis exploratorio de datos (EDA)
\end{itemize}

\vspace{1em}
\textbf{Paquetes de R a instalar}:
\begin{lstlisting}
install.packages(c("ggplot2", "dplyr", "corrplot", "GGally"))
\end{lstlisting}
\end{frame}

\end{document}
